
\subsubsection{6.1 The Inverted-U Pattern}

The finding that moderate accommodation predicts highest engagement challenges linear extrapolations from CAT. Several explanations are possible:

\begin{enumerate}
\item \textbf{Mimicry aversion}: Excessive accommodation may be perceived as artificial, potentially triggering disengagement.
\item \textbf{Optimal distinctiveness}: Users may prefer AI that adapts while maintaining some distinctiveness.
\item \textbf{Content confounds}: Formal topics may have systematically different continuation patterns independent of accommodation.
\end{enumerate}

The present study cannot distinguish among these explanations. The inverted-U pattern is correlational, and the mechanism linking accommodation to continuation remains unverified.

\subsubsection{6.2 Asymmetric Bidirectional Effects}

AI-to-human accommodation (r=0.152) substantially exceeds human-to-AI adaptation (r=0.013). This asymmetry suggests that AI systems, through training, develop stronger accommodation tendencies than humans naturally exhibit toward AI interlocutors. Designers cannot rely on natural user accommodation; the burden of adaptation falls on AI systems.

\subsubsection{6.3 Limitations}

Several limitations constrain interpretation:

\begin{enumerate}
\item \textbf{Correlational, not causal}: Association is observed; causation is not established. The study does not manipulate accommodation levels.
\end{enumerate}

\begin{enumerate}
\item \textbf{Single dataset}: Results are based on the Anthropic hh-rlhf dataset only. The hh-rlhf dataset contains chosen vs rejected response pairs, potentially introducing selection bias. Generalization to other platforms, models, or datasets is untested.
\end{enumerate}

\begin{enumerate}
\item \textbf{Engagement proxy}: Continuation is used as an engagement proxy. Continuation does not necessarily indicate satisfaction---users may continue conversations due to dissatisfaction requiring clarification.
\end{enumerate}

\begin{enumerate}
\item \textbf{Content confounds}: Formal topics may have systematically different continuation patterns independent of accommodation effects.
\end{enumerate}

\begin{enumerate}
\item \textbf{Single formality measure}: Results depend on DeBERTa-based formality scoring. Other operationalizations of formality or accommodation may yield different patterns.
\end{enumerate}

\begin{enumerate}
\item \textbf{Untested causal mechanism}: The planned hypothesis (H-M4) testing whether accommodation causally predicts conversation length was not executed due to H-M3 gate failure.
\end{enumerate}

